
In this section, we introduce our design of the  service hand-off framework based on live migration of the containers. We will first briefly introduce the different stages of service hand-off in our framework. Then in section \ref{design:mapImage}, we will discuss how we enable the sharing of image layers between two different Docker hosts with different cache IDs. In section \ref{design:memDiff}, \ref{design:compression}, and \ref{design:pipe}, we will introduce the details of how we can further speed up the migration by transferring dirty memory, compression, pipelined and paralleled processing during the migration of Docker containers.

\subsection{System Overview}


% % \subfigcapmargin = .5cm
% \begin{figure}
% \begin{subfigure}[b]{0.4\textwidth}
%     % \centering
%     \includegraphics[width=\linewidth]
%     % {figure/test-tarssh.eps}
%     {figure/sys.png}
%     \caption{Before Migration}
%     \label{subfig:before}
% \end{subfigure}
% \quad\quad\quad\quad
% \begin{subfigure}[b]{0.4\textwidth}
%     % \centering
%     \includegraphics[width=\linewidth]
%     % {figure/test-tarssh.eps}
%     {figure/sys2.png}
%     \caption{After Migration}
%     \label{subfig:after}
% \end{subfigure}
% \caption{ Offloading Serivce Handoff: Before and After Live Migration of Offloading Container. }
% \label{fig:before-after}
% \end{figure}

% https://tex.stackexchange.com/questions/350729/software-architecture-picture-in-tikz
\tikzset{
    srv/.style  = {draw, rectangle, minimum width=1.8cm, 
    minimum height=2cm, 
    text depth=2cm,
    text centered, 
    inner sep=1, 
    text width=2.0cm, font=\footnotesize,
    draw=black, 
    fill=blue!30},
    % https://tex.stackexchange.com/questions/6899/tikz-multi-line-text-in-the-node-description
    conNode/.style   = {draw, rectangle, minimum width=1.6cm, minimum height=1cm, text centered, text width=1.5cm, font=\footnotesize, draw=black,
    fill=yellow!30 
    % fill=white
    },
    dockerNode/.style = {draw, rectangle, minimum width=1.6cm, minimum height=0.5cm, text centered, text width=1.5cm, font=\footnotesize, draw=black, fill=yellow!30
    },
    cloudNode/.style = {draw,
    cloud, 
    cloud puffs=25.7, 
    cloud ignores aspect, 
    minimum width=3.5cm,
    minimum height=0.5cm, 
    align=center,
    fill=white
    },
}
\begin{figure}

\begin{subfigure}[b]{0.23\textwidth}
    %  https://tex.stackexchange.com/questions/115087/how-can-i-embed-an-external-image-within-a-tikzpicture-environment
    % https://tex.stackexchange.com/questions/91320/arrows-from-cloud-cross-the-outer-line
    % \begin{scaletikzpicturetowidth}{\textwidth}
    \begin{tikzpicture}[thick,scale=0.8, every node/.style={scale=0.8}]
    \coordinate (cloudco) at (1.6,1.5);
    \coordinate (serverAco) at (0,3.5);
    \coordinate (serverBco) at (2.5,3.5);
    \coordinate (clientco) at (0.8,0);
    % \draw[->] (cloudco) -- (90:2);
    % server node A
    \node [name=srvA, srv,draw] at (serverAco) {Edge Server A};
    \node [name=conA, conNode] at (srvA.center) {Offloading Container};
    \node [name=dockerA,dockerNode] at ([yshift=1em]srvA.south) {Docker};
    % draw line arrow from client to serverA
    \draw[->] (clientco) -- (dockerA.south);
    % server node B
    \node [name=srvB, srv] at (serverBco) {Edge Server B};
    \node [name=dockerB,dockerNode] at ([yshift=1em]srvB.south) {Docker};
    \node (phone) at (clientco)
        {\includegraphics[width=.2\textwidth]{figure/fig.jpg}};
    \node[cloudNode] (cloud) at (cloudco) {WAN};
    \draw[->] (phone) -- node[above]{moving} ++ (\textwidth/2,0);
    \end{tikzpicture} 
    \caption{Before Migration}
    \label{subfig:before}
\end{subfigure}
\begin{tikzpicture}[thick]
\draw[dashed] (0,0) -- (0,5.3);
\end{tikzpicture} 
\begin{subfigure}[b]{0.23\textwidth}
%  https://tex.stackexchange.com/questions/115087/how-can-i-embed-an-external-image-within-a-tikzpicture-environment
    % https://tex.stackexchange.com/questions/91320/arrows-from-cloud-cross-the-outer-line
    \begin{tikzpicture}[thick,scale=0.8, every node/.style={scale=0.8}]
    \coordinate (cloudco) at (1.6,1.5);
    \coordinate (serverAco) at (0,3.5);
    \coordinate (serverBco) at (2.5,3.5);
    \coordinate (clientco) at (1.5,0);
    % \draw[->] (cloudco) -- (90:2);
    % server node A
    \node [name=srvA, srv,draw] at (serverAco) {Edge Server A};
    \node [name=dockerA,dockerNode] at ([yshift=1em]srvA.south) {Docker};
    % server node B
    \node [name=srvB, srv] at (serverBco) {Edge Server B};
    \node [name=conB, conNode] at (srvB.center) {Offloading Container};
    \node [name=dockerB,dockerNode] at ([yshift=1em]srvB.south) {Docker};
    % draw line arrow from client to serverA
    \draw[->] (clientco) -- (dockerB.south);
    \node (phone) at (clientco)
        {\includegraphics[width=.2\textwidth]{figure/fig.jpg}};
    \node[cloudNode] (cloud) at (cloudco) {WAN};
    \draw[->] (phone) -- node[above]{moving} ++ (\textwidth/2,0);
    \end{tikzpicture}

\caption{After Migration}
\label{subfig:after}
\end{subfigure}
\caption{ Offloading Serivce Handoff: Before and After Migration of Offloading Container. }
\label{fig:before-after}
\end{figure}


Figure~\ref{fig:before-after} shows an example usage scenario of offloading service hand-off based on Docker container migration. Before migration, mobile client offloads its computations to its nearest edge server \textbf{A}, where all the computations are processed inside a Docker container. When the mobile client moves far away from server \textbf{A} and reaches the service area of edge server \textbf{B}, it's offloading computation will be migrated from server \textbf{A} to server \textbf{B}. This is done via migration of its offloading Docker containers, where all the runtime memory states as well as its storage data are transferred to the target server \textbf{B}.

% shows the overview of the whole migration stages. Generally, 
Figure~\ref{fig:migstages} shows the architecture in different processing stages described below:
% our architecture supports high speed offloading service hand-off through the following steps:
% \begin{enumerate*}[series = tobecont, itemjoin = \quad]

\begin{enumerate}[label=\textbf{S\arabic*}]

\item \label{predownload}\textbf{Pre-download Application Image.} Upon the container starts and offers offloading service to the client, we also dynamically predict the possible target edge servers that are near to the client. Then we request those target servers to download the Docker image of the offloading application from Docker Hub.

\item \label{predump} \textbf{Pre-dump Container.} Before we received the migration request, we will try to dump a snapshot of the container's runtime memory and send to the possible target edge servers we found in stage \ref{predownload}. The container is still running during this stage. 

\item \label{prepare} \textbf{Remapping Layer Cache IDs.} Upon receiving the hand-off request to a certain target server, prepare for file system synchronization by remapping the image layers' local identifications on the target server.

% \item \label{firstsync} Upon receiving migration request, compress the thin container layer and send to target.

\item \label{checkpoint} \textbf{Checkpoint and Stop Container.} After we remapped the layer IDs according the shared original layer IDs between source and target edge servers, we then could checkpoint and stop the container in order to send all the latest run-time states to the target edge server.
From this point the container will stop running on the source server and waiting to be restored on the target server.

\item \label{fs-sync} \textbf{Contaienr Layer Synchronization} After checkpointed the container, the file system will not be changed. In this step, we will send the container layer's contents to the target server. At the same time, all the checkpointed runtime states and configuration files, such as \textit{state.json, config.v2.json, etc.} are also transferred to the target server. 

\item \label{daemon-Reload} \textbf{Docker Daemon Reload} After we send all the runtime state and configuration files to the target node, Docker daemon on the target node still cannot recognize the container immediately.  This is because those runtime state and configuration files are created by another Docker daemon on source node, the daemon on target node doesn't have those configurations loaded into its runtime database. So we need to reload the Docker daemon in order to reload those runtime state and configuration files it just received from the source node.

\item \label{img-sync} \textbf{Get Memory Difference.} After we get checkpointed images from the final dump of the container, we then compare this final dump memory to the pre-dumped memory in stage \ref{predump}. We then could get the memory difference and send only these difference to the target server. 
% \item \label{finalsync} Synchronize the file system for the last time, updating only the changed files of the container layer before last synchronization in \ref{firstsync}.

\item \label{restore} \textbf{Restore Container.} Finally, we will restore container on the target host with all its runtime status retrieved. At the same time, the system will go to stage \ref{predownload} to prepare the next iteration of service handoff in the future.

\end{enumerate}


% (Before migration, both the source and target edge server have the application base images downloaded.)
 
\subsection{Match Shared Image Layers} \label{design:mapImage}

As we introduced above, in the AUFS-based addressable layers, the local cached ID for each downloaded image layer is stored in 

\subsection{Get Memory Difference} \label{design:memDiff}
In order to reduce the total memory image size during hand-off, we checkpoint the container and get a snapshot of the container's memory in stage \ref{predump}. The snapshot was transferred before the hand-off starts, just as the application's Docker images, which are also downloaded before the hand-off starts.

Upon the hand-off starts, we checkpoint the container and stop it. Then we do a diff operation on the new memory dump images along with the old snapshot we got from stage \ref{predump}. Then we only need to transfer the difference of the dump memory to the target node. We use Xdelta3 to get the memory difference.

\begin{figure}
\begin{subfigure}[b]{0.45\textwidth}
    % \centering
    \includegraphics[width=\linewidth]
    % {figure/test-tarssh.eps}
    {figure/face-518744dump_iter_11.eps}
    \caption{Memory Dumps}
    \label{fig:dump}
\end{subfigure}

\begin{subfigure}[b]{0.45\textwidth}
    % \centering
    \includegraphics[width=\linewidth]
    % {figure/test-tarssh.eps}
    {figure/face-518744-iter-11.eps}
    \caption{Dirty Memory, original}
    \label{fig:diff-ori}
\end{subfigure}
\begin{subfigure}[b]{0.45\textwidth}
    % \centering
    \includegraphics[width=\linewidth]
    % {figure/test-tarssh.eps}
    {figure/face-518744-iter-11-2.eps}
    \caption{Dirty Memory, adjacent}
    \label{fig:diff-adj}
\end{subfigure}

\begin{subfigure}[b]{0.45\textwidth}
    % \centering
    \includegraphics[width=\linewidth]
    % {figure/test-tarssh.eps}
    {figure/face-518744diffdiff_ori11.eps}
    \caption{Dirty Difference, original}
    \label{fig:diffdiff-adj}
\end{subfigure}
\begin{subfigure}[b]{0.45\textwidth}
    % \centering
    \includegraphics[width=\linewidth]
    % {figure/test-tarssh.eps}
    {figure/face-518744diffdiff_adj11.eps}
    \caption{Dirty Difference, adjacent}
    \label{fig:diffdiff-ori}
\end{subfigure}

\caption{  Dirty Memory Size Analysis for OpenFace. The data shows 11 dumps of the running container of OpenFace, indexed from dump 0 to dump 10. Memory is dumped every 10 seconds. 
Dirty Memory Size of Dump 1~10, Compared to Original Dump 0 .
Dirty Memory Size Between Two Adjacent Dumps. Memory is dumped every 10 seconds.
Size of Difference Between Two Adjacent Dirty Memory from Every three Adjacent Memory Dumps. Memory is dumped every 10 seconds. }
\label{fig:diff-diff}
\end{figure}

% \begin{figure}[tb!]
%     \centering
%     \includegraphics[width=0.9\textwidth]
%     % {figure/test-tarssh.eps}
%     {figure/face-518744-iter-11.eps}
%     \caption{Dirty Memory Size of Dump 1~10, Compared to Original Dump 0 }
%     \label{fig:diff-0}
% \end{figure}

% \begin{figure}[tb!]
%     \centering
%     \includegraphics[width=0.45\textwidth]
%     % {figure/test-tarssh.eps}
%     {figure/face-518744-iter-11-2.eps}
%     \caption{Dirty Memory Size Between Two Adjacent Dumps. Memory is dumped every 10 seconds }
%     \label{fig:diff-adjacent}
% \end{figure}

% \begin{figure}[tb!]
%     \centering
%     \includegraphics[width=0.45\textwidth]
%     % {figure/test-tarssh.eps}
%     {figure/face-518744-diffdiff-11.eps}
%     \caption{Size of Difference Between Two Adjacent Dirty Memory from Every three Adjacent Memory Dumps. Memory is dumped every 10 seconds. }
%     \label{fig:diff-diff}
% \end{figure}


\subsection{Compression of the Transferred Data}\label{design:compression}

During the live migration of a container, we mainly have 4 kinds of data need to transfer:

The layer stacks information, the container writable layer, the meta data files of the container, and the main memory images checkpointed by CRIU. 

% In order to speed up the synchronization process, we combine both the compression techniques and synchronization tools. 

The layer stacks information is send via RPC socket connection. Those information is only a list of SHA256 ID strings, so it's quite efficient to be sent as a socket message. There is no need to compress it.

The container writable layer and meta data files are regular files in the storage system, so we use bzip2  to compress and send via ssh connection.

The dump memory images are binary data, which is not efficient to compressed using bzip2. So we use ??? to compress and send via ssh connections.

\subsection{Pipelined \& Parallel Processing}\label{design:pipe}


% \ref{predump}, \ref{img-sync}.

% For the root file system and meta data files. We have two stages of files synchronization before and after the container was checkpoints. The first synchronization in \ref{firstsync} will transfer the base file system for the container layer, so it's better to compress it. We use 'tar' command to compress and sent via SSH. For the final synchronization of the file system, since we already have most of the files transferred, we choose to use 'rsync' to send over only the changed files.

% For the checkpointed memory images, we compress them via 'tar' compression and send via SSH connections.

% Syncing smaller files individually through rsync or scp results in each file starting at least one own data packet over the net. If the file is small and the packets are many, this results in increased protocol overhead. Now count in that there are more than one data packets for each file by means of rsync protocol as well (transferring checksums, comparing...), the protocol overhead quickly builds up

% \url{http://unix.stackexchange.com/questions/30953/tar-rsync-untar-any-speed-benefit-over-just-rsync}

% TODO: more details about compression and synchronization: more memory images compression techniques (to replace tar cmd). 

% TODO: rsync vs. tar comparison

